%% =============================================================================
%% Agency, Communion, and Proxemics as Universal Control Variables for
%% Socially Expressive Robots
%%
%% Author : Chengxu Zhou (Humanoid Robotics Lab, UCL)
%% Source : reverse-engineered from docs/proposal/social_HRI.pdf (3-page proposal)
%% Format : IEEEtran (journal mode)
%%
%% This LaTeX is a faithful transcription — content unchanged from the PDF.
%% =============================================================================

\documentclass[journal]{IEEEtran}

\usepackage{cite}
\usepackage{amsmath,amssymb,amsfonts}
\usepackage{graphicx}
\usepackage{textcomp}
\usepackage{xcolor}
\usepackage{booktabs}
\usepackage{url}
\usepackage{hyperref}

\begin{document}

\title{Agency, Communion, and Proxemics as Universal Control Variables for Socially Expressive Robots}

\author{Chengxu~Zhou \thanks{All authors are with the Humanoid Robotics Lab, Department of Computer Science, University College London, UK.}}

\markboth{Manuscript}{Zhou: Agency, Communion, and Proxemics for Socially Expressive Robots}

\maketitle

\begin{abstract}
Robots today excel at task execution yet remain poor at being \emph{understood} in social contexts. We propose a principled framework that parametrises robot actions within a three-dimensional social space defined by Agency, Communion, and Proxemics (ACP). We theoretically ground these variables in social psychology, map them onto robot motion control, and introduce a reinforcement learning and generative modelling pipeline to realise socially expressive behaviours. Through controlled human--robot interaction studies, we demonstrate that ACP variables systematically alter human perception of robots' trustworthiness, politeness, and cooperation. Our findings establish ACP as a minimal universal set of control variables for socially expressive embodied AI.
\end{abstract}

\begin{IEEEkeywords}
Humanoid robots, embodied AI, reinforcement learning, motion generation, social interaction, agency, communion, proxemics.
\end{IEEEkeywords}

%% ===========================================================================
\section{Introduction}\label{sec:intro}
\IEEEPARstart{R}{obots} are increasingly deployed in human-centered environments, yet their acceptance and trust hinge not only on functional competence but on their ability to communicate social intent~\cite{belkaid2021gaze}. Humans interpret actions along latent socio-affective dimensions; however, robot control traditionally optimises only task-level goals (\emph{e.g.}, stability, precision). This disconnect leads to socially awkward or unsafe interactions.

\textbf{Our proposition:} grounding robot motion generation in three fundamental social variables---\emph{Agency} (assertiveness), \emph{Communion} (affiliativeness), and \emph{Proxemics} (interpersonal distance)~\cite{wiggins1991agency, hall1966hidden}. We hypothesise that ACP coordinates determine how humans perceive and cooperate with robots, and that parametrising robot control in this space can yield predictable, interpretable, and safe social behaviours.

\textbf{Core novelty.} Unlike affect labels or single-cue heuristics (\emph{e.g.}, gaze), we elevate \emph{social variables to first-class control parameters} that are optimised jointly with physical feasibility. We further show that these variables constitute a compact control basis that generalises across platforms (humanoid / manipulator-on-quadruped) and tasks (passing, handover, side-by-side walking).

%% ===========================================================================
\section{Related Work}\label{sec:related}

\subsection{Embodied AI and Social Interaction}
Large language/vision-language models have recently been extended to robots, but their social grounding remains weak. \emph{Science Robotics} has highlighted that safe social interaction requires principles of interpretability and uncertainty awareness~\cite{belkaid2021gaze}.

\subsection{Social Dimensions in Psychology}
\begin{itemize}
\item \textbf{Agency--Communion}: canonical axes of interpersonal behaviour~\cite{wiggins1991agency}.
\item \textbf{Proxemics}: interpersonal distance zones (intimate, personal, social, public)~\cite{hall1966hidden}.
\item \textbf{Valence--Arousal--Dominance}: widely used in affective computing~\cite{russell1980circumplex}, but less actionable in control.
\end{itemize}

\subsection{Robotics and Social Variables}
Prior work has studied gaze~\cite{belkaid2021gaze}, interpersonal distance, or affect recognition, but no generalisable control framework has been proposed that formalises multiple social variables as \emph{first-class control parameters}.

blabla\dots

%% ===========================================================================
\section{Theoretical Framework}\label{sec:framework}
We formalise ACP as follows. Let $a \in [0,1]$, $c \in [0,1]$, and $p \in [0.5, 2.5]\,\mathrm{m}$ be continuous control variables.

\begin{itemize}
\item \textbf{Agency} ($a$): mapped to motion onset latency $\Delta t_{\text{start}}$, peak velocity $v_{\max}$, and torque allocation/actuation aggressiveness. High $a=$proactive, forceful; low $a=$hesitant, reactive.
\item \textbf{Communion} ($c$): mapped to trajectory smoothness (low jerk), lateral yielding displacement $\Delta y$, and openness of posture (arm abduction and torso yaw). High $c=$affiliative, cooperative.
\item \textbf{Proxemics} ($p$): mapped to the target interpersonal distance band (personal/social/public). Enforced as a soft constraint on robot--human separation $d(t)$ around a target band $[d^{-}, d^{+}]$ centred at $p$.
\end{itemize}

\subsection{Operationalisation via Kinematic Proxies}
Given a motion segment $X_{1:T}$ (skeleton or task-space markers), we compute weak labels:
\begin{align}
a &= \alpha_1 \left(1 - \tfrac{\Delta t_{\text{start}}}{T_0}\right)
      + \alpha_2 \,\mathrm{norm}(v_{\max})
      + \alpha_3 \,\mathrm{norm}(\bar{a}), \label{eq:a-proxy} \\
c &= \beta_1 \,\mathrm{norm}(\text{arm-open})
      + \beta_2 \left(1 - \mathrm{norm}(\text{jerk})\right) \nonumber \\
  &\quad + \beta_3 \,\mathrm{norm}(\Delta y / \Delta x), \label{eq:c-proxy} \\
p &= \gamma_1 \,\mathrm{norm}(d_{\min})
      + \gamma_2 \,\mathrm{norm}(\bar{d})
      + \gamma_3 \,\mathrm{norm}(d_{\text{stop}}), \label{eq:p-proxy}
\end{align}
where $\mathrm{norm}(\cdot)$ is per-dataset z-score then min--max scaling.

\subsubsection*{Remark}
When datasets lack explicit counterparts, $p$ can be supplied at control time while $a, c$ are supervised using proxies only.

%% ===========================================================================
\section{Technical Realisation}\label{sec:tech}

\subsection{Conditional Motion Generator}
We seek $G(\text{scene}, a, c, p) \to \text{traj}$, with two implementation lines.

\subsubsection{Line A: Conditional Regression (MVP)}
A temporal Transformer maps $(\text{scene}, a, c, p)$ to a low-dimensional reference trajectory. Loss combines trajectory reconstruction, conditional alignment, and physical feasibility.

\subsubsection{Line B: Conditional Motion Diffusion}
We adopt a DiT-style denoiser with classifier-free guidance on $(a, c, p)$, plus conditional alignment regularisation and projected guidance for physical plausibility.

\subsection{Socially-Constrained RL Control}
Low-level controllers are trained in Isaac Lab (PPO/SAC) to track generator outputs while optimising social rewards:
\begin{align}
r &= w_{\text{task}} r_{\text{task}}
   + w_{ag} r_{\text{Agency}} \nonumber \\
  &\quad + w_{com} r_{\text{Communion}}
   + w_{prox} r_{\text{Proxemics}}, \label{eq:reward} \\
r_{\text{Agency}} &= f_{ag}(\Delta t_{\text{start}}, v_{\max}),
   \label{eq:r-agency} \\
r_{\text{Communion}} &= f_{com}(\Delta y / \Delta x, \text{jerk}^{-1}),
   \label{eq:r-com} \\
r_{\text{Proxemics}} &= -\|d(t) - p\|^{2}. \label{eq:r-prox}
\end{align}

\subsection{Uncertainty-Aware Gating}
We estimate perception uncertainty via MC-dropout on the human pose estimator. If entropy exceeds threshold, fallback policy is triggered: $(a, c, p) = (a_{\min}, c_{\text{mid}}, p_{\text{public}})$.

\subsection{Cross-Platform Deployment}
We validate on humanoid and quadruped-with-arm platforms using the same $(a, c, p)$ interface, only tuning WBC parameters.

%% ===========================================================================
\section{Experiments}\label{sec:exp}

\subsection{Tasks}
\begin{enumerate}
\item Passing in a corridor (Proxemics, Communion).
\item Object handover (Agency, Proxemics).
\item Side-by-side walking (Agency, Communion).
\end{enumerate}

\subsection{Datasets and Implementation}
AMASS for base kinematics; BEHAVE and in-house MoCap for proxemics calibration. Temporal Transformer baseline; DiT diffusion upgrade. RL trained in Isaac Lab with 4096 envs.

\subsection{Metrics}
\begin{itemize}
\item \textbf{Objective:} separation, yielding displacement, completion time, stability, energy.
\item \textbf{Behavioural:} user acceptance, avoidance, pace synchrony.
\item \textbf{Subjective:} 7-point Likert on politeness, trust, cooperativeness.
\end{itemize}

\subsection{User Study Protocol}
\subsubsection{Participants}
30 participants (balanced gender, age 20--40) recruited under institutional ethics approval.
\subsubsection{Procedure}
Each participant experiences robots performing tasks under different $(a, c, p)$ conditions and baseline controllers. Trials are randomised.
\subsubsection{Measures}
Participants complete post-trial questionnaires (Likert scales). Behavioural data (acceptance, avoidance) are logged via motion capture.
\subsubsection{Analysis}
Mixed-effects ANOVA across conditions; Holm correction for multiple comparisons. Effect sizes reported.

\subsection{Baselines}
Task-only, Valence--Arousal generation, end-to-end LLM planner.

\subsection{Ablations}
Remove social rewards, disable generator, disable gating, swap ACP with alternative style conditions.

%% ===========================================================================
\section{Results}\label{sec:results}
ACP settings produce consistent shifts in perception and behaviour; high $a$, high $c$, social-zone $p$ maximise cooperation. Mismatched ACP yields avoidance. Uncertainty gating reduces missteps under occlusion. Portability confirmed across platforms.

%% ===========================================================================
\section{Discussion}\label{sec:discussion}

\subsection{Scientific Contribution}
ACP is validated as a universal social control basis: continuous, optimisable, and platform-agnostic.

\subsection{Limitations and Future Work}
Tasks are short scripted; cultural differences not studied. Proxemics supervision partly synthetic.

\subsection{Reproducibility and Ethics}
Code and data dictionaries will be released. User study under ethics approval, informed consent obtained.

%% ===========================================================================
\section{Conclusion}\label{sec:conclusion}
We establish Agency, Communion, and Proxemics as first-class variables in robot motion generation. By uniting generative modelling, socially-constrained RL, and uncertainty-aware gating, we demonstrate ACP governs human interpretation of robot behaviour.

%% ===========================================================================
\bibliographystyle{IEEEtran}
\bibliography{references}

\end{document}
